\begin{example}[7.2.1.1]
\label{IV.2.s7.2.1.1-fr}
Sea $A$ un anillo local noetheriano \emph{de dimensión} 1. Entonces
las condiciones siguientes son equivalentes:
\begin{enumerate}
\item[\emph{a)}] $A$ carece de ideales primos asociados inmersos.
\item[\emph{a$'$)}] $\widehat A$ carece de ideales primos asociados inmersos.
\item[\emph{b)}] $A$ es estrictamente equidimensional.
\item[\emph{b$'$)}] $\widehat A$ es estrictamente equidimensional.
\item[\emph{c)}] $A$ es un anillo de Cohen-Macaulay.
\end{enumerate}
\end{example}

En efecto, \emph{a)} significa que el ideal maximal $\mathfrak m$ de $A$ no
está asociado a $A$; por tanto, los ideales primos asociados a $A$ son
los ideales minimales de $A$, todos distintos de $\mathfrak m$, y
\oldpage[IV-2]{188}
para un ideal $\mathfrak p$ de esta clase se tiene necesariamente
$\dim(A/\mathfrak p)=1$; así, \emph{a)} implica \emph{b)}. Recíprocamente, \emph{b)}
implica que todo ideal primo $\mathfrak p\in\operatorname{Ass}(A)$ es
distinto de $\mathfrak m$ y, por ello, minimal; en consecuencia, \emph{b)}
implica \emph{a)}. Ya se sabe que \emph{a)} y \emph{c)} son equivalentes para un anillo local
de dimensión 1 (5.7.8). Como afirmar que $A$ es un anillo de
Cohen-Macaulay equivale a afirmar que $\widehat A$ es un anillo de Cohen-Macaulay
(\textbf{0}, 16.5.2), y se tiene $\dim(\widehat A)=1$, se ve finalmente que
\emph{a$'$)} y \emph{b$'$)} son equivalentes a \emph{c)}.
